\documentclass[conference]{IEEEtran}
\usepackage[T1]{fontenc}
\usepackage[brazil]{babel}
\usepackage{amsmath,booktabs,url,balance}
\usepackage[hidelinks]{hyperref}

\title{GraphRAG para consulta rastreável de exposições financeiras em relatórios de companhias brasileiras}
\author{\IEEEauthorblockN{Arysson André Oliveira da Luz \quad Caio Justino de Andrade}
\IEEEauthorblockN{Luiz Henrique Monteiro dos Santos \quad Paulo Victor Oliveira de Barros Araújo}
\IEEEauthorblockA{Centro Universitário de João Pessoa (Unipê)\\
Tendências em Ciência da Computação\\
Professor Ricardo Roberto de Lima}}

\begin{document}
\maketitle
\begin{abstract}
Apresentamos um protótipo de consulta a exposições financeiras que combina recuperação textual e vetorial, consultas Cypher geradas por um modelo de linguagem e caminhos em um grafo de propriedades no Neo4j. O acervo reúne seis relatórios oficiais de Suzano, Klabin e Irani, referentes a 2024 e 2025, com 954 páginas. A implementação indexa 3.918 trechos narrativos elegíveis e mantém 85 declarações com procedência. Regras preservam identidade, exercício, versão, negação e condições reconhecidas. Caminhos financeiros participam da seleção do contexto antes da geração, e uma interface permite conferir afirmações, relações e páginas da fonte. Uma verificação de recuperação com duas perguntas de desenvolvimento localizou a evidência esperada em uma pergunta na variante vetorial e em duas nas variantes híbrida e com caminhos. O algoritmo não aumentou o recall em relação ao híbrido nesse conjunto. Esses resultados demonstram operação no recorte estudado, sem comprovar superioridade, generalização ou redução de alucinações. A revisão independente e a avaliação reservada continuam pendentes.
\end{abstract}
\begin{IEEEkeywords}
GraphRAG, grafos de conhecimento, recuperação de informação, exposições financeiras, procedência.
\end{IEEEkeywords}

\section{Introdução}
Relatórios financeiros associam obrigações, moedas, indexadores e medidas de proteção a empresas, exercícios e condições específicos. Uma consulta que localiza palavras semelhantes pode recuperar passagens úteis sem preservar todas essas relações. Por exemplo, identificar um instrumento de proteção em uma nota não permite concluir que a companhia eliminou seu risco ou protegeu todos os passivos. A resposta precisa respeitar o objeto e o alcance da declaração.

Este trabalho investiga a pergunta: como combinar texto e relações financeiras para produzir respostas com evidências rastreáveis? Desenvolvemos um protótipo que transforma descrições narrativas em um grafo de propriedades e utiliza esse grafo na preparação do contexto de um modelo de linguagem. Cada afirmação retornada pode identificar fontes e caminhos inspecionáveis.

O objetivo implementado é consultar exposições descritas nos documentos. O protótipo não estima probabilidade de insolvência nem calcula um ranking de risco financeiro. A denominação inicial de mapa de riscos foi delimitada a um mapa de relações e exposições declaradas. Essa delimitação acompanha a cobertura efetivamente validada dos dados.

Este artigo é uma versão para revisão, com resultados de desenvolvimento até 15 de novembro de 2026. Descrevemos a arquitetura e as verificações disponíveis, distinguindo sua conclusão operacional das etapas científicas ainda necessárias. O código e os registros são mantidos no repositório privado do grupo\cite{repositorio}.

\section{Fundamentação e posicionamento}
Edge et al.\cite{edge} apresentam uma abordagem GraphRAG voltada a perguntas globais sobre coleções textuais. Seu método constrói um grafo de entidades, organiza comunidades e utiliza resumos dessas comunidades na geração. Nosso protótipo adota um esquema financeiro e caminhos dirigidos como mecanismo de enriquecimento do contexto. Portanto, não reproduz o pipeline de sumarização por comunidades desse trabalho.

Liang et al.\cite{liang} propõem KAG para combinar modelos de linguagem e grafos em domínios profissionais, incluindo representação de conhecimento, ligação entre grafo e trechos e raciocínio híbrido. Nosso projeto se aproxima desses princípios ao explicitar conceitos e regras financeiras e ao manter as relações ligadas às fontes. Não implementamos o framework KAG completo ou seu mecanismo de formas lógicas.

O interesse acadêmico do protótipo está na integração verificável entre extração, busca, algoritmo e resposta. A existência de um grafo, isoladamente, não garante melhoria. O experimento precisa distinguir localização de evidência, suporte de cada afirmação e adequação da abstenção.

\section{Acervo e limites da informação}
Selecionamos demonstrações financeiras completas em português de três companhias da cadeia de papel, celulose e embalagens\cite{suzano,klabin,irani}. Suzano e Klabin constam no cadastro consultado da CVM como Papel e Celulose; Irani consta como Embalagens. Essa diferença impede tratar a seleção como um único subsetor homogêneo. A Tabela~\ref{tab:corpus} apresenta as páginas dos PDFs efetivamente coletados.

\begin{table}[t]
\caption{Acervo documental coletado}
\label{tab:corpus}
\centering
\begin{tabular}{lrrr}
\toprule
Empresa & 2024 & 2025 & Total\\
\midrule
Suzano &184&181&365\\
Klabin &153&156&309\\
Irani &139&141&280\\
\midrule
Total &476&478&954\\
\bottomrule
\end{tabular}
\end{table}

Um manifesto versionado registra CNPJ, código CVM, exercício, protocolo DFP, versão declarada, URL oficial e SHA-256 de cada PDF. O protocolo é uma conferência de emissor, exercício e versão; o hash identifica o arquivo publicado pelo emissor que foi usado, inclusive quando esse arquivo reúne material adicional de divulgação. As fontes e os protocolos podem ser conferidos no manifesto e nos dados abertos da CVM\cite{cvm}.

Os arquivos contêm 3.682.930 caracteres extraídos. Das 954 páginas, 19 têm menos de cem caracteres. Esses indicadores auxiliam a triagem, mas não demonstram integridade de tabelas. A disposição de certas tabelas produz valores concatenados no texto, e uma figura pode conter informação relevante sem texto extraível.

O recorte atual usa 3.918 trechos classificados como narrativa elegível. A quantidade de páginas coletadas não equivale à cobertura semântica do sistema. Tabelas, imagens, unidades e colunas comparativas ainda exigem validação específica. Por isso, a aplicação restringe respostas quantitativas e não calcula exposição líquida. O extrator também precisa de validação em documento externo ao conjunto usado no desenvolvimento.

\section{Método e implementação}
\subsection{Extração e procedência}
O pipeline verifica o checksum, extrai o texto com pypdf, preserva páginas e posições e segmenta o documento. Regras de domínio reconhecem entidades e relações financeiras, produzindo declarações com sujeito, predicado, objeto, companhia, período, escopo, polaridade, condição e evidência quando identificados. NER e RE são automáticos, porém dependem de vocabulário e construções previstos nas regras.

Uma entidade representa um conceito identificável; uma declaração representa uma observação documentada. O modelo separa \texttt{Empresa}, \texttt{Documento}, \texttt{Trecho} e \texttt{Declaracao} de classes financeiras como \texttt{ObrigacaoFinanceira}, \texttt{MedidaProtecao}, \texttt{Moeda} e \texttt{FatorRisco}. Uma aresta financeira guarda o ID da declaração que a sustenta. A declaração aponta para o trecho por \texttt{EVIDENCIADA\_POR}, e o trecho aponta para seu documento por \texttt{PERTENCE\_A}.

O reconhecimento de menções e relações foi verificado em oito excertos com quinze menções e vinte e três relações esperadas, anotados com assistência de IA. Esses excertos participaram da correção das regras. Portanto, seus resultados são de desenvolvimento, e não uma medida independente de precisão ou generalização. Dois formulários de revisão humana foram preparados, mas não foram preenchidos até esta versão.

\subsection{Identidade, versão e escopo}
A companhia usa CNPJ normalizado como identidade. Moedas, indexadores e fatores podem ser conceitos globais; obrigações e medidas de proteção são vinculadas à companhia. Constraints de unicidade impedem IDs repetidos, mas não substituem a avaliação semântica das decisões de união e separação de menções.

Documento, empresa, tipo e exercício definem uma série versionada. A publicação de uma revisão preserva o histórico e retira a anterior da busca corrente na mesma transação. Reprocessar uma versão antiga não a reativa. Testes verificaram substituição, idempotência e rollback com revisões sintéticas. As seis fontes reais continuam na versão declarada um.

O escopo tem uma limitação relevante: as 85 declarações atuais estão qualificadas como não especificado pelo extrator. A regra procura qualificadores explícitos na sentença e não recupera suficientemente o contexto de cabeçalhos e notas. Esse resultado descreve o reconhecimento automático; não significa ausência de escopo nos relatórios. Para evitar combinações indevidas, caminhos com escopo desconhecido usam declarações evidenciadas no mesmo trecho e limitam a confiança da resposta.

\subsection{Banco e recuperação híbrida}
Utilizamos Neo4j 5.26 Community, com labels, relacionamentos tipados, propriedades, constraints e índices Range, Fulltext e vetorial. Os vetores têm 384 dimensões e usam cosseno. A implementação de embeddings utiliza MiniLM multilíngue em ONNX quantizado, média com máscara, normalização L2 e truncamento em 256 tokens. Os artefatos e revisões estão fixados no catálogo do projeto.

A recuperação aplica filtros de empresa, exercício e vigência. Busca vetorial e Fulltext consultam o mesmo universo narrativo elegível. Uma consulta Cypher completa também é produzida pelo LLM a partir da pergunta e do esquema permitido. Essa consulta passa por contrato de leitura, parâmetros, limite de resultados e inspeção do plano com \texttt{EXPLAIN}; execução e reparo ficam no rastreio. O banco opera em modo somente leitura durante as consultas.

Combinamos as posições dos candidatos por Reciprocal Rank Fusion:
\begin{equation}
\mathrm{RRF}(t)=\sum_{m\in M:t\in L_m}\frac{1}{60+r_m(t)},
\end{equation}
em que $L_m$ é a lista do método $m$ e $r_m(t)$ é a posição do trecho. Cada sinal retorna até doze candidatos, e a fusão seleciona até oito trechos. A fusão por posição evita somar diretamente scores de escalas diferentes. O limite de 256 tokens do encoder pode truncar partes relevantes de trechos longos e constitui uma limitação de representação.

\subsection{Algoritmo e formação do contexto}
Aplicamos \texttt{shortestPath} sobre relações financeiras dirigidas permitidas, com até quatro saltos. Há caminhos de dois saltos verificados no acervo. A seleção verifica empresa, exercício, vigência, escopo e as declarações que sustentam cada aresta. Relações de procedência e de substituição documental não contam como saltos financeiros.

O contexto inclui fontes e caminhos antes de o modelo gerar a resposta. Removemos duplicatas e prefixos cobertos por caminhos maiores com as mesmas arestas e o mesmo escopo. A distribuição por empresa e exercício reduz a concentração das fontes de uma única companhia em perguntas comparativas. O orçamento total, incluindo relações, é de 3.500 caracteres.

Um fator canônico comum pode aparecer em declarações de empresas diferentes. Isso permite inspecionar um conceito compartilhado; não estabelece transação entre elas, risco equivalente, correlação ou contágio. O sistema proíbe transformar esse compartilhamento em uma conclusão financeira sem evidência adicional.

\subsection{Geração e interface}
O modelo local é Qwen2.5-1.5B-Instruct em GGUF Q4\_K\_M, servido por llama.cpp b7770 em CPU, contexto de 4.096 tokens, temperatura zero e seed 42. A saída estruturada contém até três afirmações, fontes, caminhos, cobertura e possível conflito. O catálogo preserva revisões e checksums dos modelos.

Guardas determinísticas verificam referências, suporte literal de números, algumas conclusões absolutas, relevância por vocabulário e qualificação de escopo. A aplicação preserva a saída original e todas as transformações. As categorias sem suporte, limitada e moderada são definidas por regras e não representam probabilidades calibradas. Esses controles não verificam toda implicação semântica.

A interface Streamlit apresenta resposta, confiança, fontes, Cypher e um subgrafo interativo das relações efetivamente consultadas. A seleção de uma afirmação realça seu caminho e permite abrir a página do PDF. Tarefas continuam após desconexão do navegador em um deployment de processo único; a reinicialização do servidor interrompe a inferência. Esse comportamento foi verificado localmente, sem teste de carga.

\section{Protocolo de avaliação}
Comparamos três variantes: vetor, híbrido e híbrido com caminhos. O orçamento de contexto é igual. A variante com caminhos reutiliza os candidatos híbridos, permitindo observar o enriquecimento do contexto sem repetir uma geração Cypher que poderia mudar os candidatos.

Na recuperação, uma evidência esperada é localizada quando há correspondência de documento, página e termo normalizado no texto. O recall é a fração de localizadores esperados encontrados, medida nos candidatos e nas fontes efetivamente enviadas. Um localizador representa uma evidência conhecida, não o conjunto exaustivo de passagens que poderiam apoiar a resposta.

O smoke contém duas perguntas já usadas no desenvolvimento, ambas da Irani em 2025: instrumentos de proteção de juros (S01) e ligação de ACC com recebíveis de exportações (S02). A rodada registra hashes do conjunto, corpus, código e relatório, modelos, versões e consultas geradas. Uma tentativa anterior com timeout foi preservada como inválida para comparação.

A avaliação de respostas usa os contextos congelados da recuperação e as mesmas regras de geração e validação. Falhas técnicas são registradas separadamente de abstenções. Suporte semântico, cobertura, correção de empresa e exercício e adequação da abstenção dependem de revisão. Não calculamos automaticamente uma taxa de alucinação a partir de referências válidas.

Para a rodada científica, o protocolo prevê quinze perguntas de desenvolvimento e trinta reservadas, com distribuição entre consulta direta, multi-hop, comparação, temporalidade, negação e falta de suporte. O conjunto, o gabarito e a configuração precisam ser congelados antes das execuções finais, com revisão independente e registro das divergências.

\section{Resultados de desenvolvimento}
\subsection{Recuperação com indexação completa}
A rodada com 3.918 vetores concluiu as seis execuções sem falhas. As duas consultas Cypher foram aprovadas na primeira tentativa. A Tabela~\ref{tab:recuperacao} apresenta localização da evidência esperada no contexto. Cada pergunta tem um localizador esperado.

\begin{table}[t]
\caption{Evidência esperada no contexto do smoke}
\label{tab:recuperacao}
\centering
\begin{tabular}{lccc}
\toprule
Pergunta & Vetor & Híbrido & Com caminhos\\
\midrule
S01 &0/1&1/1&1/1\\
S02 &1/1&1/1&1/1\\
\bottomrule
\end{tabular}
\end{table}

O algoritmo acrescentou um caminho em cada pergunta, sem aumento de recall sobre o híbrido. Todos os contextos ficaram abaixo de 3.500 caracteres. O conjunto é pequeno e conhecido, de uma única companhia, e não permite inferir superioridade ou generalização.

\subsection{Seleção entre empresas e operação da interface}
Em um diagnóstico anterior, com o índice piloto de 553 passagens e os mesmos candidatos e arestas, a distribuição do contexto por empresa e a compactação de prefixos aumentaram de duas para três as companhias representadas. O contexto passou de 3.444 para 3.391 caracteres. Esse caso motivou uma correção de implementação; não é comparação independente entre métodos nem resultado da indexação ampliada.

Uma consulta real sobre ACC e recebíveis da Irani apresentou uma resposta parcial, confiança limitada e caminho de dois saltos com fonte na página PDF 124 de 2025. A consulta foi retomada após recarregamento do navegador sem repetir a inferência. Sua latência foi 335,269 segundos no computador local, com disponibilidade restrita de memória. Trata-se de um caso operacional anterior à expansão do índice, não de um benchmark controlado.

\subsection{Erros de contrato na geração}
Uma rodada inicial sobre os seis contextos congelados concluiu a inferência, mas quatro respostas das variantes sem grafo foram rejeitadas por referências a caminhos inexistentes. O modelo confundiu IDs de fontes ou trechos com IDs de caminhos. As duas respostas com grafo passaram pelas verificações determinísticas, o que não demonstra que todo seu conteúdo esteja correto.

O diagnóstico levou a um contrato de decodificação contextual: o campo de fontes aceita apenas IDs de fontes disponíveis, e o de caminhos aceita apenas caminhos existentes ou uma lista vazia quando não houver caminhos. A rodada original foi preservada e a comparação posterior usa esse contrato igualmente nas três variantes. Nenhuma conclusão de redução de alucinação semântica é extraída desse ajuste.

A segunda rodada concluiu seis inferências sem referências de caminhos inexistentes. Cinco respostas finais mantiveram afirmações e cobertura parcial; a variante com grafo em S01 foi rejeitada porque a saída declarava cobertura insuficiente junto de afirmações. O contrato corrigiu o tipo de referência observado, mas esse novo resultado também mostra que o grafo não garante uma resposta final aproveitável. Os seis itens foram preparados para revisão de suporte documental, ainda pendente.

\section{Limitações e ameaças à validade}
A seleção de três companhias e dois exercícios limita a validade externa. As perguntas do smoke são conhecidas e concentram-se em uma empresa. Os excertos anotados participaram do desenvolvimento, e a revisão humana permanece pendente. Esses fatores impedem apresentar métricas de generalização de NER, RE ou resposta.

O extrator baseado em regras pode omitir relações, interpretar objetos de forma incompleta e perder qualificadores distantes da sentença. A falta de escopo reconhecido nas 85 declarações restringe as travessias. A narrativa extraída também pode conter resíduos de gráficos e tabelas. Quantidade de vetores e ausência de erros de execução não demonstram qualidade financeira das informações.

O modelo local pequeno pode copiar passagens irrelevantes, produzir relações inadequadas e declarar cobertura excessiva. Restringir IDs impede certas referências inválidas, mas não comprova apoio semântico de uma afirmação em uma fonte existente. O pós-processamento pode ainda reduzir cobertura ao remover afirmações ou transformar uma resposta em abstenção.

A comparação de latência exige rodadas controladas: aquecimento, cache e disputa de memória variam. O tempo adicional de enriquecimento por grafo não equivale ao tempo total da variante. As sementes e artefatos fixados favorecem rastreabilidade, sem garantir reprodução bit a bit em qualquer hardware.

\balance
\section{Conclusão e continuidade}
O protótipo integra extração narrativa, Neo4j, recuperação híbrida, Cypher gerado, algoritmo de caminhos e uma interface que permite inspecionar relações e fontes. As verificações demonstram funcionamento no acervo estudado e participação do grafo no contexto enviado ao modelo.

O smoke não demonstrou aumento de recall do algoritmo sobre a recuperação híbrida. A análise de geração identificou um erro de contrato que exigiu correção comum às variantes. A contribuição científica final depende da validação em documento externo, do tratamento de tabelas e escopo, da revisão independente e da execução reservada. A conclusão será atualizada de acordo com esses resultados, inclusive se não houver benefício do grafo em determinadas classes de perguntas.

\section*{Reprodutibilidade e assistência de IA}
O repositório contém configuração Docker, dependências fixadas do ambiente Windows, fontes e hashes, scripts de ingestão, indexação, consulta e avaliação, testes e instruções. PDFs integrais, pesos, credenciais e execuções locais ficam fora do Git. As fontes podem ser readquiridas e verificadas pelo manifesto. O acesso do professor ao repositório privado depende de autorização posterior do grupo.

A implementação e a redação receberam assistência de IA. As anotações de desenvolvimento são identificadas como assistidas. A revisão acadêmica dos autores, sua participação prática, a validação independente e a defesa presencial não são substituídas por essa assistência.

\begin{thebibliography}{9}
\bibitem{edge} D. Edge et al., ``From Local to Global: A Graph RAG Approach to Query-Focused Summarization,'' arXiv:2404.16130, 2024. [Online]. Disponível: \url{https://arxiv.org/abs/2404.16130}.
\bibitem{liang} L. Liang et al., ``KAG: Boosting LLMs in Professional Domains via Knowledge Augmented Generation,'' arXiv:2409.13731, 2024. [Online]. Disponível: \url{https://arxiv.org/abs/2409.13731}.
\bibitem{cvm} Comissão de Valores Mobiliários, ``Demonstrações Financeiras Padronizadas e cadastro de companhias abertas,'' Portal de Dados Abertos. [Online]. Disponível: \url{https://dados.cvm.gov.br/}. Acesso: 14 nov. 2026. URLs por documento e checksums no manifesto do projeto.
\bibitem{repositorio} A. A. O. da Luz, C. J. de Andrade, L. H. M. dos Santos e P. V. O. de B. Araújo, ``GraphRAG riscos financeiros,'' repositório privado, 2026. [Online]. Disponível sob autorização: \url{https://github.com/victort01/graphrag-riscos-financeiros}. Registros: \texttt{data/manifesto.json}, \texttt{data/modelos.json} e resumos de avaliação versionados.
\bibitem{suzano} Suzano S.A., ``Demonstrações financeiras completas 2024 e 2025,'' Central de Resultados. [Online]. Disponível: \url{https://ir.suzano.com.br/Portuguese/informacoes-financeiras/central-de-resultados/default.aspx}. PDFs e hashes no manifesto.
\bibitem{klabin} Klabin S.A., ``Demonstrações financeiras 2024 e 2025,'' Central de Resultados. [Online]. Disponível: \url{https://ri.klabin.com.br/divulgacoes-e-resultados/central-de-resultados/}. PDFs e hashes no manifesto.
\bibitem{irani} Irani Papel e Embalagem S.A., ``Demonstrações financeiras 2024 e 2025,'' Relações com Investidores. [Online]. Disponível: \url{https://ri.irani.com.br/}. PDFs e hashes no manifesto.
\end{thebibliography}
\end{document}
